% Promoted after complete independent root and companion proof audits.
\subsection{Generic sharpness for primitive prime polarizations}

Work over an algebraically closed field \(k\) of characteristic two.
Global generation is an explicit hypothesis throughout this subsection.
Relative Frobenius on the surface and absolute Frobenius on the
coefficient space are denoted \(F_A\) and \(F_X\), respectively.

\begin{lemma}[Smooth members of an odd globally generated polarization]
\label{lem:primitive-odd-smooth-members}
\status{Proved; independent audits passed.}
Let \(L\) be an ample globally generated line bundle on an abelian
surface \(A\), with odd \(N=h^0(A,L)>1\). The complete system
\(|L|\) has a nonempty open of smooth connected members, of genus
\(N+1\).
\end{lemma}
\begin{proof}
The odd theta kernel has a maximal isotropic subgroup of order
\(N\). Its theta extension splits, and the resulting actual
linearization descends \(L\) along an etale isogeny
\(\eta:A\to B\) of degree \(N\) to an ample line bundle \(M\).
Since \(\chi(M)=\chi(L)/N=1\), the descended polarization is
principal. These are the odd theta-descent facts used in
Lemma~\ref{lem:DE-primitive-one-three-smooth-member}; they depend
on the oddness of the kernel, not on its having order three.

The complete principal-theta alternative proved in that lemma
gives either a smooth genus-two theta divisor or two smooth
elliptic translates meeting transversely at one point. Its etale
pullback is therefore a reduced divisor \(C_0\in|L|\) that is
smooth or nodal. This assertion permits disconnected inverse
images of individual components and does not impose a nonproduct
condition on \(B\).

If there are nodes, choose \(s_1\in H^0(A,L)\) nonzero at every
one of them. Global generation and the infinitude of \(k\) allow
avoidance of these finitely many proper linear hyperplanes.
For \(C_0=(s_0=0)\), consider the pencil \(s_0+z s_1\).
At a central node its completed equation is
\(xy+z u(x,y,z)=0\), with \(u\) a unit. On the relative
singular scheme the partial derivatives put \(x,y\) in \((z)\);
substitution in the equation then puts a unit times \(z\) in
its ideal. Thus no relative singular component at that node
dominates the parameter trait. Away from the nodes the central
fibre is smooth, and properness excludes any other generic
singular point. The generic pencil member is smooth. If \(C_0\)
was already smooth, openness gives the same conclusion directly.

The hypersurface sequence and ample-line vanishing
\(H^1(A,L^{-1})=0\) give \(H^0(C,\mathcal O_C)=k\) for
every effective member. Hence a smooth member is connected.
Adjunction and \(L^2=2N\) give genus \(N+1\).
This proof uses explicit smoothing, not an unqualified
positive-characteristic Bertini assertion.
\end{proof}

\begin{lemma}[The whole-space Euler map and its kernel criterion]
\label{lem:primitive-prime-Euler-criterion}
\status{Proved; independent audits passed.}
Let \(A\) be nonordinary, let \(L\) have type \((1,p)\) for
an odd prime \(p\), and put \(W=H^0(A,L)\),
\(X=\mathbf P(W)\). The quotient of the universal divisor's
coherent \(H^1\) by the constant ambient part is
\(\mathcal E=\Omega_X^1\), with a regular linearized Frobenius
\(\Phi:F_X^*\mathcal E\to\mathcal E\).
In a level-\(p\) theta basis, its chartwise formula is
\[
 \Phi_s=M_s^t\big|_{(s^{[2]})^\perp}:
             (s^{[2]})^\perp\longrightarrow s^\perp,
 \qquad M_s(v)=\beta(s,v).                         \tag{PGR1}
\]
If \(M_s\) has rank \(p-1\) and
\(\ker M_s^t=k y_s\), then
\[
       \Phi_s\text{ is bijective}\quad\Longleftrightarrow\quad
                    \sum_i y_{s,i}s_i^2\ne0.       \tag{PGR2}
\]
These assertions are valid for every nonzero coefficient section,
without requiring its divisor to be smooth.
\end{lemma}
\begin{proof}
The universal hypersurface sequence, ample-line vanishing and
Serre duality identify the quotient as
\[
 \ker\bigl(W^*\otimes\mathcal O_X(-1)
                      \longrightarrow\mathcal O_X\bigr)
                    =\Omega_X^1.                  \tag{PGR3}
\]
Frobenius preserves the constant ambient subbundle and is regular
on this kernel. Duality with relative Cartier of products gives
(PGR1), exactly as in the complete whole-space hypersurface
calculation of Lemma~\ref{lem:DE-primitive-one-three-Euler}.
That calculation uses the actual Frobenius pullback of the Euler
frame and not the numerical dimension three. The source is the
annihilator of \(s^{[2]}\) in \((W^{(2)})^*\), whereas the
target is the annihilator of \(s\) in \(W^*\).

The tensor \(\beta:\bigwedge^2W\to W^{(2)}\) is nonzero
by Theorem~\ref{thm:DE-odd-polarization-Cartier-surjection}.
Its actual theta actions have central weights one and two, so
Lemma~\ref{lem:prime-Heisenberg-two-contractions} gives exactly
the matrices (PHR2). In particular \(M_s s=0\), so the image
of \(M_s^t\) is contained in \(s^\perp\).
When its rank is \(p-1\), restriction between the two displayed
spaces of dimension \(p-1\) is injective precisely when its
one-dimensional kernel is not contained in \((s^{[2]})^\perp\).
This is exactly (PGR2).
\end{proof}

\begin{proposition}[Birational radicals give generic Euler bijectivity]
\label{prop:primitive-prime-generic-Euler}
\status{Proved; independent audits passed.}
Under the preceding geometric hypotheses, suppose the actual
dual alternating pencil \(B_y\) has generic rank \(p-1\).
Then \(\Phi\) is generically bijective on \(X\).
This conclusion holds unconditionally under those geometric
hypotheses for
\[
                         p\in\{3,5,7,11,13,17\}.   \tag{PGR4}
\]
\end{proposition}
\begin{proof}
Let \(\sigma(y)=(\sigma_i(y))\) be the nonzero maximal-Pfaffian
vector of \(B_y\). Each entry has degree \((p-1)/2\).
Theorem~\ref{thm:prime-Heisenberg-birational-radical} proves that
\([y]\mapsto[\sigma(y)]\) is birational, and that \(M_s\)
has generic rank \(p-1\). There are therefore dense opens on
which this map and its inverse are defined and the two ranks
are maximal. On them \(s=\sigma(y)\) represents the section
and \(y\) represents \(\ker M_s^t\).

Consider the homogeneous polynomial of degree \(p\)
\[
                       Q(y)=\sum_i y_i\sigma_i(y)^2.
                                                        \tag{PGR5}
\]
In characteristic two its ordinary derivatives are
\[
                           \frac{\partial Q}{\partial y_i}
                                    =\sigma_i(y)^2.  \tag{PGR6}
\]
Since \(\sigma\) is not the zero vector, \(Q\) is not the
zero polynomial. Its nonvanishing locus meets the dense common
birational open. At every point of that intersection (PGR2)
proves that the actual Euler Frobenius is bijective. Hence its
determinant is not the zero section and its bijective locus is
a nonempty open of \(X\).

The universal support certificates of
Proposition~\ref{prop:prime-Heisenberg-dual-rank-through-seventeen}
supply the assumed dual rank for the listed primes, because the
actual Cartier tensor was proved nonzero in the preceding lemma.
No inference from a finite set of coefficient values is used.
\end{proof}

\begin{theorem}[Generic curve rank for the certified primitive primes]
\label{thm:primitive-prime-generic-curve-rank}
\status{Proved; independent audits passed.}
Let \(A\) be a fixed nonordinary abelian surface in characteristic
two and \(L\) an ample globally generated line bundle of type
\((1,p)\), where \(p\) is in (PGR4). On a nonempty open of
smooth connected members \(C\in|L|\),
\[
                   g(C)=p+1,\qquad
                   f(C)=f(A)+p-1.                  \tag{PGR7}
\]
The same conclusion holds for any odd prime for which the actual
dual Cartier pencil has generic rank \(p-1\).
\end{theorem}
\begin{proof}
The smooth-member open is nonempty by
Lemma~\ref{lem:primitive-odd-smooth-members}; the Euler-bijective
open is nonempty by the proposition. They meet because \(X\)
is integral. On their intersection the actual coherent sequence
\[
 0\longrightarrow H^1(A,\mathcal O_A)
 \longrightarrow H^1(C,\mathcal O_C)
 \longrightarrow\mathcal E_{[C]}\longrightarrow0
                                                        \tag{PGR8}
\]
is Frobenius compatible, and its final term is bijective of
dimension \(p-1\). The stable fixed-vector functor is exact
over the algebraically closed field: on the nilpotent part the
additive lifting operator is inverted by a finite geometric sum,
and on the bijective part the additive Lang map is surjective.
Thus the dimensions of stable parts add in (PGR8), giving
\(f(C)=f(A)+p-1\). Adjunction gives \(g(C)=p+1\).
\end{proof}

\begin{remark}[Scope of generic bijectivity]
\label{rem:primitive-prime-generic-rank-scope}
\status{Proved; independent audits passed.}
The statement retains global generation and the actual nonordinary
Cartier tensor. It does not assert global generation for every
line bundle of type \((1,p)\), and does not prove a generic-rank
statement for an untested prime without its stated dual-rank
hypothesis. In particular a source of type \((1,17)\) is not
identified with the complete principal system \(|51\Theta|\).

Only the nonvanishing open of (PGR5) is used. Although (PGR6)
makes its zero divisor smooth on the dual maximal-rank open,
this proof does not show that this divisor meets the birational
isomorphism open: it might have a component contracted by the
radical map. No common maximal-rank wall, full local radical,
nonzero-Lang connectedness, or full monodromy group is inferred.
No new numerical experiment is needed. The complete previously
audited theta-descent, primitive Euler, odd-Cartier surjection and
prime-Heisenberg proofs are the dependencies.
\end{remark}
